% ID: SEL-D0024-OB-E04
% ESTADO: validada. S025, cotejada con examen y criterios.
\textbf{a)} Los diez valores suman 65 horas, así que $\bar x=6.5$.
Como la población es normal y $\sigma=2$ es conocida, usamos
$\bar x\pm z_{1-\alpha/2}\sigma/\sqrt n$. Al $90\%$,
$z_{0.95}\approx1.644854$. Resulta
\[
6.5\pm 1.644854\frac{2}{\sqrt{10}}
\approx\boxed{[5.460,\ 7.540]\text{ horas}}.
\]
\textbf{b)} Para un margen de como máximo $0.1$ horas,
$z_{0.95}2/\sqrt n\le0.1$. Esto equivale a
$n\ge(20z_{0.95})^2\approx1082.2174$. El primer entero
que cumple la condición es $\boxed{1083\text{ alumnos}}$.
